\section{Entire initial generation of the same exceptional family}

Let $R=\mathbb C[[x_1,\ldots,x_n]]$, with its actual coordinate maximal
ideal $\mathfrak m$, and let $\rho_e$ be the literal projection setting the
exceptional coordinate $x_e$ equal to zero. Homogeneous components and all
ideal spans below use the whole original formal series rings.

\begin{theorem}[Exceptional filtered division generates the entire initial ideal]
\label{hir:actual-exceptional-family-entire-initial-generation}
\uses{hir:actual-exceptional-saturated-homogeneous-fibre-filtered-division,hir:actual-homogeneous-initial-component-derivative-dictionaries}
\lean{Hironaka.HironakaActualExceptionalInitialGeneration.actualHomogeneousInitialIdeal_eq_exceptional_component_span}
\leanok
For a finite family $(H_i)$ with assigned degrees $d_i$, suppose
$H_i\in\mathfrak m^{d_i}$, each $\rho_e H_i$ is homogeneous of degree
$d_i$, and $x_eF\in(H_i)$ implies $F\in(H_i)$ for every whole $F$.
Then
\[
 \operatorname{in}_{\mathfrak m}(H_i)
   =\bigl((H_i)_{d_i}:i\bigr).
\]
The left side is the entire homogeneous initial ideal of the actual whole
span, and the right side uses the degree components of these same generators.
\end{theorem}
\begin{proof}
The exceptional filtered division theorem gives, for every
$F\in(H_i)\cap\mathfrak m^q$, an exact whole identity
$F=\sum_iH_iQ_i$ with $Q_i\in\mathfrak m^{q-d_i}$. If $d_i\le q$,
the product homogeneous-component formula gives
$(H_iQ_i)_q=(H_i)_{d_i}(Q_i)_{q-d_i}$, since all lower coefficients in
both factors vanish. If $d_i>q$, the whole product has order at least
$d_i>q$, so its degree-$q$ component is zero. Taking the component of the
finite identity puts $F_q$ in the displayed component span. For every
nonzero $F\in(H_i)$, choose its first nonzero degree $q$; its homogeneous
initial form is this same $F_q$. These forms generate the whole initial
ideal, proving the forward inclusion. Conversely each
$H_i\in(H_i)\cap\mathfrak m^{d_i}$, so its degree-$d_i$ component
belongs to the whole homogeneous initial ideal by actual degree lifting.
Taking the whole span proves the reverse inclusion, including generators
whose degree component is zero and empty families.
\end{proof}

\begin{theorem}[The same current forms have genuine homogeneous multiplier bands]
\label{hir:actual-exceptional-family-homogeneous-multiplier-bands}
\uses{hir:actual-exceptional-family-entire-initial-generation,hir:actual-homogeneous-family-shifted-filtered-division}
\lean{Hironaka.HironakaActualExceptionalInitialGeneration.exists_exceptional_initial_homogeneous_bands}
\leanok
Under the same hypotheses, if $P$ is homogeneous of degree $q$ and belongs
to the whole current initial ideal, then
\[
 P=\sum_i(H_i)_{d_i}A_i,
\]
where $A_i$ is homogeneous of degree $q-d_i$ and $A_i=0$ whenever
$d_i>q$. Thus the coefficients are actual finite homogeneous polynomials;
no polynomial relation or chosen finite representation is supplied.
\end{theorem}
\begin{proof}
The preceding equality puts $P$ in the whole span of the homogeneous
$(H_i)_{d_i}$. Homogeneous filtered division gives a whole representation
$P=\sum_i(H_i)_{d_i}B_i$ with
$B_i\in\mathfrak m^{q-d_i}$. Take its degree-$q$ component. If $d_i\le q$,
homogeneity of the first factor leaves exactly
$(H_i)_{d_i}(B_i)_{q-d_i}$. If $d_i>q$, every coefficient of the product
at degree $q$ is zero. Set $A_i=(B_i)_{q-d_i}$ in the first case and
$A_i=0$ in the second. These are homogeneous of the stated degrees, and
homogeneous series in finitely many variables are their literal finite
polynomials. The resulting finite identity is the assertion.
\end{proof}

\begin{theorem}[Original canonical point quotients give the actual entire initial span]
\label{hir:actual-original-point-canonical-current-initial-generation}
\uses{hir:actual-canonical-point-normalized-strict-filtered-division,hir:actual-exceptional-family-entire-initial-generation}
\lean{Hironaka.HironakaActualExceptionalInitialGeneration.exists_original_point_canonical_initial_family}
\leanok
An original analytic germ ideal $J$ and an actual point chart choose one
normalized canonical family $H_a$, before all current order bounds. Its
whole span is the independently defined normalized-member ideal. Whenever
each $H_a$ retains its own original vertex degree $d_a$, the entire
homogeneous initial ideal of its actual Taylor extension equals
$((\widehat H_a)_{d_a}:a)$. The family, original ideal and chart are the
same in the normalization equations, whole span and initial equality.
\end{theorem}
\begin{proof}
Choose the original canonical point family supplied by actual exceptional
filtered division. Its whole analytic span is the independently defined
normalized ideal, and the theorem supplies all shifted formal representations
for this same family under exactly the actual current order bounds. The
component argument in the first proof applies to these representations and
the Taylor order bounds. It identifies every actual initial form with the
same component span; its reverse inclusion follows from actual generator
membership and order. No second normalized family is chosen.
\end{proof}

\begin{theorem}[Literal relative quotients give the same point initial generation]
\label{hir:actual-relative-point-same-family-current-initial-generation}
\uses{hir:actual-original-point-canonical-current-initial-generation,hir:actual-relative-point-prescribed-normalized-family}
\lean{Hironaka.HironakaActualExceptionalInitialGeneration.canonical_relative_point_initial_generation}
\leanok
Fix the original canonical analytic generators $G_a$ of $J$, an actual
relative chart with selected coordinates $A$, pivot $e$ and direction $b$,
and prescribed actual quotients $e^{d_a}H_a=\beta_{A,e,b}^*G_a$.
Let $\sigma$ scale precisely the tangent coordinates outside $A$ by $e$
and fix all the normal coordinates. If every $\sigma H_a$ has actual
maximal order at least $d_a$, then the entire homogeneous initial ideal of
$(\widehat{\sigma H_a})$ is
$((\widehat{\sigma H_a})_{d_a}:a)$.
\end{theorem}
\begin{proof}
The literal chart composition gives
$e^{d_a}\sigma H_a=\beta_{\mathrm{point},e,b_A}^*G_a$, with the original
prescribed direction on $A$ and zero direction outside $A$. Choose the
canonical point family from the preceding theorem for exactly that chart.
Both families satisfy these same equations with the same numerator $G_a$.
The actual analytic germ ring is a domain and the pivot germ is nonzero;
cancelling $e^{d_a}$ therefore identifies each point quotient with
$\sigma H_a$. Transfer the genuine current order bounds and the initial
equality along this pointwise equality. This proves the assertion for the
prescribed relative family, rather than for an independently chosen witness.
\end{proof}

These are statements about the actual current family and its whole initial
ideal. The finite image-witness construction, total filtered descent to the
relative chart, and the complete current-centre Hilbert--Samuel criterion
require further arguments.
